%% HAND-WRITTEN circuitikz figure -- NOT generated by maestro2tex. The tool only places
%% it, via the "order" list in the block's maestro2tex config; the `sch_` prefix keeps it
%% clear of the generated fig_*/tab_* namespace.
%%
%% Depicts: the signal-path topology of the class-AB transimpedance amplifier
%% castalia/anatop_tia_amp_ab (34 transistors, copied from teewinot/AdcBufferOpamp) --
%% complementary NMOS+PMOS input pairs, the two folded-cascode branches, the class-AB
%% output pair M15/M14 driven from ABP/ABN, the cascode (Ahuja) compensation nodes C1
%% and C2, and the two external cc_ab capacitors that close the compensation from C1 and
%% C2 onto Out.
%%
%% Topology transcribed on 2026-08-16 from subckt anatop_tia_amp_ab in
%%   ic/.simulation/castalia/tb_anatop_tia_amp_ab/maestro/results/maestro/
%%     Interactive.14/1/TiaFbTB/netlist/input.scs   (Spectre order M(drain gate source bulk)).
%% The same file's top level gives the external compensation:
%%   OPAMP (net1 net2 vdd! WE V_CM V_OUT ...) anatop_tia_amp_ab   -> C1=net1, C2=net2,
%%   C2 (net1 V_OUT) capacitor c=cc_ab ; C3 (net2 V_OUT) capacitor c=cc_ab ; cc_ab=500f.
%%
%% Drawn: 22 of 34 transistors. Omitted: the BiasSwitchCascode instance and the four
%% power-down devices M57/M58/M48/M49, and the class-AB bias generator M50--M53, M55,
%% M56, M59, M60 that develops the DiodeN/DiodeP references. Internal netlist node names
%% net014/net019/net026 are kept as drawn; C1, C2, ABP, ABN, InP, InM, Out, DiodeN and
%% DiodeP are real terminal/net names. BP/BPC/BN/BNC are the enable-gated forms of the
%% bp!/bpc!/bn!/bnc! globals, generated inside BiasSwitchCascode.
%%
%% Identically labelled nets are connected; wires crossing without a filled dot are not.
\begin{figure}[htbp]
  \centering
  \resizebox{\linewidth}{!}{%
\ctikzset{bipoles/length=0.9cm, transistors/scale=0.82, capacitors/scale=0.9}
\begin{circuitikz}[
    every node/.append style={font=\footnotesize},
    gatelab/.style={font=\scriptsize, inner sep=1.5pt, fill=white},
    netlab/.style={font=\scriptsize, inner sep=1.5pt, fill=white},
    devlab/.style={font=\scriptsize, text=black!60, inner sep=1pt},
    ilab/.style={font=\scriptsize, align=center, inner sep=2pt, text=black!60},
    stagelab/.style={font=\Large, align=center, inner sep=2pt, text=black!60},
    railnode/.style={circle, fill=black, inner sep=1.05pt},
    grp/.style={draw=black!45, dashed, line width=0.6pt, rounded corners=4pt},
  ]

% ---- geometry -------------------------------------------------------------
\def\xTP{-3.0}  \def\xTN{-1.9}          % tail devices M6 / M0
\def\xM{0}      \def\xP{3.2}            % InM-side / InP-side input columns
\def\xC{6.4}    \def\xG{7.9}            % mirror branch, mirror gate bus
\def\xD{9.6}    \def\xAn{8.8} \def\xAp{10.4}  % output cascodes, class-AB pair
\def\xtap{11.6} \def\xgP{12.5} \def\xgN{12.95}
\def\xO{14.3}   \def\xbox{15.5} \def\xcc{16.5} \def\xend{18.6}

\def\yvdd{15.2} \def\yvss{-1.2}
\def\yPm{13.6}  \def\ynA{12.4} \def\yCa{11.5} \def\yPc{10.3}
\def\yInN{8.6}  \def\ytn{7.5}   \def\ytp{6.5}  \def\yInP{5.4}
\def\yNc{3.7}   \def\yCb{2.5}   \def\ynB{1.6} \def\yNm{0.4}
\def\yABP{9.4}  \def\yABN{4.6}  \def\yAB{7.0}  \def\ynS{8.2}
\def\yout{7.0}  \def\ycct{14.5} \def\yccb{-0.5}


% ---- supply rails ---------------------------------------------------------
\draw[thick] (-5.0,\yvdd) -- (\xend,\yvdd) node[right]{\texttt{inh\_vdd}};
\draw[thick] (-5.0,\yvss) -- (\xend,\yvss) node[right]{\texttt{inh\_vss}};

% ===== tail current sources ===============================================
\node[pmos](M6) at (\xTP,\yPm){};
\draw (M6.source) -- (\xTP,\yvdd);
\draw (M6.drain)  -- (\xTP,\ytp) -- (\xP,\ytp);
\draw (M6.gate) -- ++(-0.45,0) node[gatelab, left]{\texttt{BP}};
\node[devlab, anchor=west] at (\xTP+0.5,\yPm){M6};

\node[nmos](M0) at (\xTN,\yNm){};
\draw (M0.source) -- (\xTN,\yvss);
\draw (M0.drain)  -- (\xTN,\ytn) -- (\xP,\ytn);
\draw (M0.gate) -- ++(-0.45,0) node[gatelab, left]{\texttt{BN}};
\node[devlab, anchor=west] at (\xTN+0.5,\yNm){M0};

% ===== InM column: M4 / M1 (nch in) | M1L (pch in) / M3L ==================
\node[pmos](M4)  at (\xM,\yPc){};
\node[nmos](M1)  at (\xM,\yInN){};
\node[pmos](M1L) at (\xM,\yInP){};
\node[nmos](M3L) at (\xM,\yNc){};
\draw (M4.source) -- (\xM,\ynA);
\draw (M4.drain)  -- (M1.drain);
\draw (M1.source) -- (\xM,\ytn);
\draw (M1L.source)-- (\xM,\ytp);
\draw (M1L.drain) -- (M3L.drain);
\draw (M3L.source)-- (\xM,\ynB);
\draw (M4.gate)  -- ++(-0.45,0) node[gatelab, left]{\texttt{BPC}};
\draw (M1.gate)  -- ++(-0.45,0) node[gatelab, left]{\texttt{InM}};
\draw (M1L.gate) -- ++(-0.45,0) node[gatelab, left]{\texttt{InM}};
\draw (M3L.gate) -- ++(-0.45,0) node[gatelab, left]{\texttt{BNC}};
\node[devlab, anchor=west] at (\xM+0.5,\yPc){M4};
\node[devlab, anchor=west] at (\xM+0.5,\yInN){M1};
\node[devlab, anchor=west] at (\xM+0.5,\yInP){M1L};
\node[devlab, anchor=west] at (\xM+0.5,\yNc){M3L};

% ===== InP column: M3 / M2 (nch in) | M1R (pch in) / M3R =================
\node[pmos, xscale=-1](M3)  at (\xP,\yPc){};
\node[nmos, xscale=-1](M2)  at (\xP,\yInN){};
\node[pmos, xscale=-1](M1R) at (\xP,\yInP){};
\node[nmos, xscale=-1](M3R) at (\xP,\yNc){};
\draw (M3.source) -- (\xP,\yCa);
\draw (M3.drain)  -- (M2.drain);
\draw (M2.source) -- (\xP,\ytn);
\draw (M1R.source)-- (\xP,\ytp);
\draw (M1R.drain) -- (M3R.drain);
\draw (M3R.source)-- (\xP,\yCb);
\draw (M3.gate)  -- ++(0.45,0) node[gatelab, right]{\texttt{BPC}};
\draw (M2.gate)  -- ++(0.45,0) node[gatelab, right]{\texttt{InP}};
\draw (M1R.gate) -- ++(0.45,0) node[gatelab, right]{\texttt{InP}};
\draw (M3R.gate) -- ++(0.45,0) node[gatelab, right]{\texttt{BNC}};
\node[devlab, anchor=east] at (\xP-0.5,\yPc){M3};
\node[devlab, anchor=east] at (\xP-0.5,\yInN){M2};
\node[devlab, anchor=east] at (\xP-0.5,\yInP){M1R};
\node[devlab, anchor=east] at (\xP-0.5,\yNc){M3R};

% tail nodes
\draw (\xM,\ytn) -- (\xP,\ytn);
\draw (\xM,\ytp) -- (\xP,\ytp);
\node[railnode] at (\xM,\ytn){};  \node[railnode] at (\xM,\ytp){};

% ===== mirror branch: M12 / M11 / M8 / M5, summing node net026 ===========
\node[pmos, xscale=-1](M12) at (\xC,\yPm){};
\node[pmos](M11) at (\xC,\yPc){};
\node[nmos](M8)  at (\xC,\yNc){};
\node[nmos](M5)  at (\xC,\yNm){};
\draw (M12.source) -- (\xC,\yvdd);
\draw (M12.drain)  -- (\xC,\ynA);
\draw (M11.source) -- (\xC,\ynA);
\draw (M11.drain)  -- (\xC,\ynS);
\draw (M8.drain)   -- (\xC,\ynS);
\draw (M8.source)  -- (\xC,\ynB);
\draw (M5.drain)   -- (\xC,\ynB);
\draw (M5.source)  -- (\xC,\yvss);
\draw (M11.gate) -- ++(-0.45,0) node[gatelab, left]{\texttt{BPC}};
\draw (M8.gate)  -- ++(-0.45,0) node[gatelab, left]{\texttt{BNC}};
\draw (M5.gate)  -- ++(-0.45,0) node[gatelab, left]{\texttt{BN}};
\node[devlab, anchor=east] at (\xC-0.55,\yPm){M12};
\node[devlab, anchor=west] at (\xC+0.5,\yPc){M11};
\node[devlab, anchor=west] at (\xC+0.5,\yNc){M8};
\node[devlab, anchor=west] at (\xC+0.5,\yNm){M5};

% net014 / net019 horizontals into the InM cascodes
\draw (\xM,\ynA) -- (\xC,\ynA);
\draw (\xM,\ynB) -- (\xC,\ynB);
\node[railnode] at (\xC,\ynA){};  \node[railnode] at (\xC,\ynB){};
\node[netlab, anchor=south] at (1.6,\ynA+0.08){\texttt{net014}};
\node[netlab, anchor=north] at (1.6,\ynB-0.08){\texttt{net019}};

% mirror gate bus: net026 -> M12.gate and M13.gate
\draw (\xC,\ynS) -- (\xG,\ynS) -- (\xG,\yPm) -- (M12.gate);
\node[railnode] at (\xC,\ynS){};  \node[railnode] at (\xG,\yPm){};
\node[netlab, anchor=east] at (\xC-0.15,\ynS){\texttt{net026}};

% ===== output cascode branch: M13 / M10 / M9 / M7 ========================
\node[pmos](M13) at (\xD,\yPm){};
\node[pmos](M10) at (\xD,\yPc){};
\node[nmos](M9)  at (\xD,\yNc){};
\node[nmos](M7)  at (\xD,\yNm){};
\draw (\xG,\yPm) -- (M13.gate);
\draw (M13.source) -- (\xD,\yvdd);
\draw (M13.drain)  -- (\xD,\yCa);
\draw (M10.source) -- (\xD,\yCa);
\draw (M10.drain)  -- (\xD,\yABP);
\draw (M9.drain)   -- (\xD,\yABN);
\draw (M9.source)  -- (\xD,\yCb);
\draw (M7.drain)   -- (\xD,\yCb);
\draw (M7.source)  -- (\xD,\yvss);
\draw (M10.gate) -- ++(-0.45,0) node[gatelab, left]{\texttt{BPC}};
\draw (M9.gate)  -- ++(-0.45,0) node[gatelab, left]{\texttt{BNC}};
\draw (M7.gate)  -- ++(-0.45,0) node[gatelab, left]{\texttt{BN}};
\node[devlab, anchor=west] at (\xD+0.5,\yPm){M13};
\node[devlab, anchor=west] at (\xD+0.5,\yPc){M10};
\node[devlab, anchor=west] at (\xD+0.5,\yNc){M9};
\node[devlab, anchor=west] at (\xD+0.5,\yNm){M7};

% C1 / C2 horizontals into the InP cascodes and out to the tap
\draw (\xP,\yCa) -- (\xtap,\yCa);
\draw (\xP,\yCb) -- (\xtap,\yCb);
\node[railnode] at (\xD,\yCa){};  \node[railnode] at (\xD,\yCb){};

% ===== class-AB floating pair ===========================================
\node[nmos](M61) at (\xAn,\yAB){};
\node[pmos, xscale=-1](M54) at (\xAp,\yAB){};
\draw (M61.drain)  -- (\xAn,\yABP);
\draw (M61.source) -- (\xAn,\yABN);
\draw (M54.source) -- (\xAp,\yABP);
\draw (M54.drain)  -- (\xAp,\yABN);
\draw (M61.gate) -- ++(-0.45,0) node[gatelab, left]{\texttt{DiodeN}};
\draw (M54.gate) -- ++(0.45,0) node[gatelab, right]{\texttt{DiodeP}};
\node[devlab, anchor=south east] at (\xAn-0.42,\yAB+0.25){M61};
\node[devlab, anchor=south west] at (\xAp+0.42,\yAB+0.25){M54};

\draw (\xAn,\yABP) -- (\xgP,\yABP);
\draw (\xAn,\yABN) -- (\xgN,\yABN);
\node[railnode] at (\xD,\yABP){};  \node[railnode] at (\xAp,\yABP){};
\node[railnode] at (\xD,\yABN){};  \node[railnode] at (\xAp,\yABN){};
\node[netlab, anchor=south] at (\xD-0.55,\yABP+0.1){\texttt{ABP}};
\node[netlab, anchor=north] at (\xD-0.55,\yABN-0.1){\texttt{ABN}};

% ===== class-AB output stage ============================================
\node[pmos](M15) at (\xO,\yPm){};
\node[nmos](M14) at (\xO,\yNm){};
\draw (M15.source) -- (\xO,\yvdd);
\draw (M14.source) -- (\xO,\yvss);
\draw (M15.drain) -- (M14.drain);
\draw (\xgP,\yABP) -- (\xgP,\yPm) -- (M15.gate);
\draw (\xgN,\yABN) -- (\xgN,\yNm) -- (M14.gate);
\node[devlab, anchor=west] at (\xO+0.5,\yPm){M15};
\node[devlab, anchor=west] at (\xO+0.5,\yNm){M14};

\draw (\xO,\yout) -- (\xend-0.6,\yout) node[right]{\texttt{Out}};
\node[railnode] at (\xO,\yout){};

% ===== external cascode-compensation capacitors =========================
\draw (\xtap,\yCa) -- (\xtap,\ycct) -- (\xcc,\ycct);
\draw (\xtap,\yCb) -- (\xtap,\yccb) -- (\xcc,\yccb);
\draw (\xcc,\ycct) to[C] (\xcc,\yout);
\draw (\xcc,\yccb) to[C] (\xcc,\yout);
\node[railnode] at (\xcc,\yout){};
\node[netlab, anchor=south] at (13.3,\ycct+0.08){\texttt{C1}};
\node[netlab, anchor=north] at (13.3,\yccb-0.08){\texttt{C2}};
\node[ilab, anchor=west, text=black!75] at (\xcc+0.35,10.75){\texttt{cc\_ab}\\\SI{500}{\femto\farad}};
\node[ilab, anchor=west, text=black!75] at (\xcc+0.35,3.25){\texttt{cc\_ab}\\\SI{500}{\femto\farad}};

% ===== group annotations ================================================
\node[stagelab, anchor=north] at (0.4,\yvss-1.1) {complementary\\input pairs and\\their cascodes};
\node[stagelab, anchor=north] at (6.4,\yvss-1.1) {diode-connected\\summing branch};
\node[stagelab, anchor=north] at (10.1,\yvss-1.1) {class-AB\\drive};
\node[stagelab, anchor=north] at (14.0,\yvss-1.1) {output};
\node[stagelab, anchor=north] at (17.1,\yvss-1.1) {testbench\\compensation};

\end{circuitikz}
  }
  \caption[{Signal-path topology of the class-AB amplifier \texttt{anatop\_tia\_amp\_ab}.}]{Signal-path topology of the class-AB amplifier \texttt{anatop\_tia\_amp\_ab}, drawn from its netlist. The complementary pairs M1/M2 (NMOS, tail M0) and M1L/M1R (PMOS, tail M6) fold into two cascode branches; the \texttt{InM} side terminates in the diode-connected branch M12--M11--M8--M5, whose summing node \texttt{net026} sets M13, which restates that current into the \texttt{InP} branch, and the floating pair M54/M61 holds \texttt{ABP} above \texttt{ABN} so both output devices M15 and M14 are actively driven. \texttt{C1} and \texttt{C2} are the common source nodes of the PMOS cascodes M10/M3 and of the NMOS cascodes M9/M3R, brought out as pins; the two \texttt{cc\_ab} capacitors that return them to \texttt{Out} are testbench components outside the cell, \SI{500}{\femto\farad} per side. Simplified: the enable switch and its four power-down devices, and the bias generator that develops \texttt{DiodeN}/\texttt{DiodeP}, are omitted, so 22 of the cell's 34 transistors are drawn; \texttt{BP}/\texttt{BPC}/\texttt{BN}/\texttt{BNC} are the enable-gated forms of the \texttt{bp!}/\texttt{bpc!}/\texttt{bn!}/\texttt{bnc!} globals and device dimensions are not shown. Identically labelled nets are connected; wires crossing without a filled dot are not.}
  \label{fig:tiaampab-topology}
\end{figure}
